\documentclass{article}

%Estos paquetes te permiten realizar correcciones del lenguaje español
\usepackage[utf8]{inputenc}
\usepackage[spanish]{babel}


%Este paquete permite alinear como vos quieras el documento
\usepackage{ragged2e}

%Estos paquetes sirven para cargar imagenes como tambien para crear subfiguras
\usepackage{graphicx}
\usepackage{caption}
\usepackage{subcaption}

%Estos dos paquetes sirven para generar bibliografia
\usepackage{biblatex} %Imports biblatex package
\usepackage{csquotes}
\addbibresource{bibliografia.bib} %Aqui se indica en que archivo LaTeX DEBE buscar las referencias bibliograficas

%Te permite trabajar con hojas en formato paisaje en vez del formato retrato 
\usepackage{pdflscape}

%Estos paquetes te permiten añadir símbolos matemáticos
\usepackage{amsmath}
\usepackage{amssymb}

%Este paquete te permite ir a una referencia con un click
\usepackage{hyperref}

%Esto permite que el formato del documento sea A4
\usepackage{geometry}
 \geometry{
 a4paper,
 total={170mm,257mm},
 left=20mm,
 top=20mm,
 }


\begin{document}

\begin{titlepage}%Inicio de la carátula, solo modificar los datos necesarios

\newcommand{\HRule}{\rule{\linewidth}{0.5mm}} 

\center 
 
%----------------------------------------------------------------------------------------
%	ENCABEZADO
%----------------------------------------------------------------------------------------

\textsc{\LARGE Universidad Nacional de Asunción}\\[1.5cm] 
\textsc{\LARGE Facultad de Ingeniería}\\[1.5cm] 

\includegraphics[width=4.5cm]{Imagenes/logo.png}\\[1cm] % Si esta plantilla se copio correctamente, va a llevar la imagen del logo de la facultad.OBS: Es necesario incluir el paquete: graphicx 
\textsc{\Large Nombre de la Cátedra}\\[1cm] % Aqui se escribe el nombre de la catedra


%----------------------------------------------------------------------------------------
%	SECCION DEL TITULO
%----------------------------------------------------------------------------------------

\HRule \\[0.4cm]
{ \huge \bfseries Aqui va el titulo}\\[0.4cm] % Titulo del documento
\HRule \\[1.5cm]
 
%----------------------------------------------------------------------------------------
%	SECCION DEL AUTOR
%----------------------------------------------------------------------------------------

\begin{minipage}{0.4\textwidth}
\begin{flushleft} \large
\emph{Profesor:}\\
%\emph{Profesores:}\\ % En caso de que sean dos profesores, eliminar la linea anterior y el primer % de esta linea
Ing. Nombre \textsc{Apellido}\\ % Nombre del Profesor
%Ing. Nombre \textsc{Apellido}\\ % En caso de tener dos profesores, quitar el primer % para que esto sea parte del codigo
\end{flushleft}
\begin{flushleft} \large
\emph{Alumno:}\\
Nombre \textsc{Apellido}\\ % Tu nombre
%Nombre \textsc{Apellido}\\ % Nombre del segundo integrante del grupo
%Nombre \textsc{Apellido}\\ % Nombre del tercer integrante del

\end{flushleft}

\end{minipage}\\[1cm]


%----------------------------------------------------------------------------------------
%	SECCION DE LA FECHA
%----------------------------------------------------------------------------------------

{\large \today}\\[2cm] % El comando \today coloca la fecha del dia, y esto se actualiza con cada compilacion, en caso de querer tener una fecha estatica, reemplazar el \today por la fecha deseada

\vfill 

\end{titlepage}

\pagebreak
%Primeramente, se crea un indice
\tableofcontents
\pagebreak
\section{Primera Sección}
Esta sección es provisoria, el editor de este documento puede cambiarla por el nombre que sea necesario. Si usted principiante o busca simplificar algunos procesos, es conveniente que seleccione \emph{plantillas.tex} y compile ese documento. Una vez compilado ese documento, revisen la sección \textbf{Introducción} para mas información.



\pagebreak

\pagebreak


\printbibliography[title={Bibliografía},heading=bibintoc]




\end{document}
